\documentclass[11pt]{article}

\usepackage[margin=1in]{geometry}
\usepackage{booktabs}
\usepackage{hyperref}
\usepackage{amsmath}
\usepackage{xcolor}
\usepackage{enumitem}

\hypersetup{colorlinks=true, linkcolor=blue, urlcolor=blue}

\title{\textbf{Cache-Aware Memory Compilation for Token-Efficient Coding Agents}}
\author{memory.cpp}
\date{Draft --- 2026}

\begin{document}
\maketitle

\begin{abstract}
Long-horizon coding agents repeatedly send large, partially stale, and
hard-to-cache repository context into the model. The cost is not only token
billing: it appears as repeated prefill work, larger attention state, lower
prompt-cache reuse, more memory pressure per request, and less reliable agent
behavior when obsolete docs, duplicate traces, or prior failed attempts enter the
prompt. We introduce \emph{Cache-Aware Memory Compilation} (CAMC), a local-first
compiler that converts repository artifacts into a typed intermediate
representation (MemoryIR), verifies freshness, filters \emph{negative memory},
canonicalizes a byte-stable prompt prefix, and selects a minimal sufficient
context under token and KV budgets. On a controlled synthetic benchmark
(CAMC-Bench) implemented entirely by the production code paths, full CAMC reduces
prompt tokens by 58.3\% versus raw context while blocking all stale,
contradictory, and secret material and preserving task success, where raw, BM25,
and embedding retrieval baselines leak unsafe context and fail the task oracle.
We are explicit that CAMC reduces \emph{prompt-side} waste and \emph{estimated} KV
pressure; it does not compress a provider's KV cache, and runtime latency must be
measured separately.
\end{abstract}

\section{Introduction}

Coding agents such as Claude Code, Codex, and Cursor are naturally
long-horizon. Across many turns they accumulate tool traces, file summaries, test
output, branch-specific details, and repeated user intent. Provider prompt caching
helps when a repeated prefix stays byte-stable, but coding-agent context routinely
mixes durable repository knowledge with volatile logs, timestamps, and task text,
defeating reuse.

The limit being hit is not only the context window; it is \emph{unmanaged context}.
A prompt can be technically valid while still being wasteful, stale, unsafe, and
difficult to cache. We argue the right abstraction is a \textbf{compiler}, not a
retrieval system: a stage that decides what belongs in the stable prefix, what
belongs in the fresh suffix, what must be excluded, and what should only be
available by rehydration handle.

\paragraph{Contributions.}
\begin{enumerate}[leftmargin=*]
  \item \textbf{MemoryIR}, a typed intermediate representation for repository
        memory carrying source anchors, trust, freshness, sensitivity, volatility,
        token cost, and cache role (\S\ref{sec:ir}).
  \item \textbf{Semantic prefix canonicalization}: a deterministic, task-independent
        stable prefix with a blake3 content hash, so the same repo state yields the
        same cache key across tasks (\S\ref{sec:prefix}).
  \item \textbf{Negative memory} and \textbf{freshness proofs}: explicit exclusion
        of stale/harmful context and source-anchored validity checks
        (\S\ref{sec:safety}).
  \item A \textbf{KV-aware optimizer} that selects context by utility-per-cost under
        token and KV budgets (\S\ref{sec:opt}).
  \item \textbf{CAMC-Bench}, a reproducible benchmark with an ablation that isolates
        the contribution of each component (\S\ref{sec:bench}).
\end{enumerate}

\section{The systems problem}

Table~\ref{tab:failures} summarizes the failure modes CAMC targets and why they
matter to an inference provider.

\begin{table}[h]
\centering
\small
\begin{tabular}{@{}p{2.6cm}p{5.2cm}p{5.2cm}@{}}
\toprule
\textbf{Failure mode} & \textbf{What happens in practice} & \textbf{Why it matters} \\
\midrule
Repeated prefill & Same repo facts and traces reprocessed each related turn. & More compute before the first token; higher latency. \\
KV-cache growth & Every unnecessary input token contributes attention state. & Memory pressure reduces concurrency; harder scheduling. \\
Cache misses & Stable facts mixed with task text, timestamps, logs, tool output. & Prompt caching and prefix reuse stop working. \\
Stale context & Obsolete docs retrieved because they look relevant. & Wrong patches, repeat loops, lower trust. \\
Unmanaged memory & Memory grows without proof, lifecycle, or contradiction handling. & Persistent memory becomes a reliability and security risk. \\
\bottomrule
\end{tabular}
\caption{Context failure modes addressed by CAMC.}
\label{tab:failures}
\end{table}

\section{MemoryIR}
\label{sec:ir}

CAMC lowers every retrieved memory into a typed \texttt{IrBlock}. Each block
records its kind (fact, invariant, decision, bug, fix, API contract, file/trace
summary, test recipe, negative memory, security warning, \ldots), a source anchor
(file, line range, commit, content hash), scalar scores in $[0,1]$ for relevance,
trust, freshness, and sensitivity, a \emph{volatility} class, a \emph{cache role}
(stable prefix, fresh suffix, or excluded), and an estimated token cost.
Representing context as typed IR --- rather than opaque retrieved strings --- is
what lets every later pass reason uniformly about cost, safety, and cache behavior,
and makes each compiled pack explainable rather than a loose bundle of text.

\section{Semantic prefix canonicalization}
\label{sec:prefix}

The \textbf{stable prefix} may contain only deterministic, low-volatility material:
repo identity, rules, durable architecture facts, API contracts, invariants,
high-confidence decisions, safe tool-use rules, and negative-memory exclusions. It
must never contain task text, current errors, failing test output, tool traces,
timestamps, or random identifiers; those constitute the \textbf{fresh suffix}.

The prefix is rendered deterministically (sections sorted, de-duplicated, whitespace
normalized) and hashed with blake3:
\[
\texttt{stable\_prefix\_hash} = \texttt{memprefix:v1:blake3:}\,H(\text{canonical prefix bytes}).
\]
Because $H$ is computed over prefix \emph{bytes} --- not the provider or task ---
the same repository state yields an identical cache key across different tasks. This
is the property that makes provider prompt caching and runtime prefix reuse actually
hit. A cache manifest records the cacheable region $[0, t_s)$ and the
non-cacheable region $[t_s, t_s + t_f)$, the repo state hash, and estimated repeated
tokens avoided under reuse.

\section{Negative memory and freshness proofs}
\label{sec:safety}

\paragraph{Negative memory.} Most systems model only what to include. CAMC records
what \emph{not} to send: stale docs, superseded decisions, wrong fixes, generated
files, duplicate context, misleading traces, security risks, prompt-injection
sources, low-signal files, and obsolete APIs. Matching candidates are excluded
deterministically with a stated reason.

\paragraph{Freshness proofs.} Before a memory enters a pack it must prove validity:
its source file and line still exist, its span hash matches, its commit is
reachable, it has not expired, no newer memory contradicts it, its source is
reliable enough, and (optionally) a human confirmed it. Each proof is serializable
with a verdict in \{fresh, stale, unverifiable\} and a proof hash, so stale or
unsupported memory is rejected pre-inference.

\paragraph{Write firewall.} Persistent memory is an attack surface. A write firewall
screens candidate writes for injection markers, secret-like material, and untrusted
low-confidence sources, quarantining them before they can be silently re-injected
into future prompts.

\section{KV-aware optimization}
\label{sec:opt}

Selection is not relevance-only. Each block is scored by utility per cost:
\[
\text{utility} = \text{relevance} \times \text{trust} \times \text{freshness} \times \text{importance},
\quad
\text{cost} = \text{tokens} \times (1 + \text{volatility}),
\quad
\text{score} = \frac{\text{utility}}{\text{cost}}.
\]
Objectives (\texttt{max\_quality}, \texttt{min\_tokens}, \texttt{kv\_efficiency},
\texttt{cache\_hit\_rate}, \texttt{freshness\_first}, \texttt{security\_first},
\texttt{minimal\_sufficient}) reweight utility and cost; compiler-style levels
\texttt{-O0..-Oz} map to objectives. Invariants and security warnings are always
retained. The optimizer greedily selects highest-scoring blocks under a token budget
and an optional KV-byte ceiling derived from a model profile.

\paragraph{KV estimate.} For a model with $L$ layers, $h$ KV heads, head dimension
$d$, and $b$ bytes per element, the KV footprint of a sequence of length $n$ is
estimated as $\text{kv\_bytes} \approx n \cdot L \cdot h \cdot d \cdot 2 \cdot b
\cdot q$, where $q$ is an optional quantization factor. We stress this is a
prompt-side \emph{estimate}: CAMC reduces the $n$ that a runtime must process; it
does not compress the provider's KV cache.

\section{CAMC-Bench}
\label{sec:bench}

CAMC-Bench compares context-selection strategies on a bundled synthetic
repo-memory corpus. The corpus contains durable facts (an invariant, a decision, a
package-manager rule, a fix, a test recipe), superseded/contradictory memories, a
secret, and distractors. Three tasks each declare an oracle: the selection must
include the task's essential memories \emph{and} include no contradictions. Every
metric in Table~\ref{tab:bench} is produced by the production IR, negative-memory,
freshness, optimizer, and KV code paths --- not hard-coded --- and is reproducible
via \texttt{memory bench camc}.

\begin{table}[h]
\centering
\small
\begin{tabular}{@{}lrrrrrrr@{}}
\toprule
\textbf{Strategy} & \textbf{Tokens} & \textbf{Cacheable} & \textbf{Reduct.} & \textbf{Stale} & \textbf{Contra.} & \textbf{Secrets} & \textbf{Success} \\
\midrule
raw full context & 638 & 0 & 0.0\% & 9 & 9 & 3 & 0/3 \\
grep & 233 & 0 & 63.5\% & 3 & 3 & 2 & 0/3 \\
BM25 & 638 & 0 & 0.0\% & 9 & 9 & 3 & 0/3 \\
embedding retrieval & 638 & 0 & 0.0\% & 9 & 9 & 3 & 0/3 \\
CAMC $-$ no negative memory & 521 & 318 & 18.3\% & 0 & 0 & 3 & 3/3 \\
CAMC $-$ no freshness proofs & 383 & 186 & 40.0\% & 9 & 9 & 0 & 0/3 \\
CAMC $-$ no canonical prefix & 266 & 0 & 58.3\% & 0 & 0 & 0 & 3/3 \\
\textbf{full CAMC} & \textbf{266} & \textbf{186} & \textbf{58.3\%} & \textbf{0} & \textbf{0} & \textbf{0} & \textbf{3/3} \\
\bottomrule
\end{tabular}
\caption{CAMC-Bench (synthetic corpus, llama-3.1-8b profile, budget 1500). Tokens:
total context tokens. Cacheable: stable-prefix tokens. Reduct.: reduction vs raw.
Stale / Contra. / Secrets: unsafe blocks present. Success: tasks passing the oracle.}
\label{tab:bench}
\end{table}

\paragraph{Reading the ablation.} Raw, BM25, and embedding baselines fit the small
corpus within budget and therefore reduce nothing, while leaking 9 stale memories,
9 contradictions, and 3 secrets and \emph{failing every task} (the oracle rejects
contradictions). Removing freshness proofs reintroduces stale and contradictory
context (success drops to 0/3); removing the canonical prefix forfeits all cacheable
tokens; removing negative memory leaks secrets. Only \textbf{full CAMC}
simultaneously minimizes tokens, preserves a cacheable prefix, blocks all unsafe
material, and passes every task.

\section{Reproducibility: ledgers and lockfiles}

Every compiled pack can emit a provenance ledger entry (included and excluded
memory with reasons, proofs passed/failed, cache regions, token tradeoffs, and a
pack hash) and a deterministic lockfile (compiler version, repo state hash, memory
ids, source and proof hashes, negative-memory version, stable-prefix and
fresh-suffix hashes, token counts). Because the stable prefix is task-independent,
\texttt{reproduce} verifies it matches byte-for-byte across runs.

\section{Related work}

\textbf{Agent memory systems} (Supermemory, Mem0, Honcho) optimize storage and
retrieval across sessions. \textbf{Retrieval} (BM25, embeddings, RAG) ranks and
returns top-$k$ candidates. \textbf{Provider prompt caching} reuses a stable prefix
at runtime. CAMC is complementary and sits \emph{after retrieval but before
inference}: it compiles, excludes, proves, and cache-optimizes the minimal
sufficient context, and it \emph{produces} the byte-stable prefix that prompt
caching needs.

\section{Limitations}

CAMC-Bench is synthetic and small by design: it measures the compilation mechanism,
not real-world win rates. Token counts use estimators unless a provider tokenizer is
enabled. KV figures are prompt-side estimates; runtime latency and memory must be
measured against a real stack (vLLM, SGLang, llama.cpp, Ollama). Freshness and
negative-memory heuristics are conservative and can over- or under-exclude; human
review remains available for high-impact memory.

\section{Conclusion}

Long-horizon coding agents do not primarily need a bigger memory store; they need a
\emph{context compiler}. By compiling repository history into freshness-proven,
security-filtered, cache-stable, minimal context --- and by measuring the result ---
CAMC reduces prompt-side waste and estimated KV pressure while preserving task
success, and leaves behind a reproducible, auditable record of exactly what the
model was asked to process.

\end{document}
